% Generic educational note template for youtube_note_generator.
% Adapt title, metadata, boxes, theorem environments, and packages to the blueprint.
\documentclass[11pt]{article}

\usepackage[utf8]{inputenc}
\usepackage[T1]{fontenc}
\usepackage{lmodern}
\usepackage[a4paper,margin=1in]{geometry}
\usepackage{microtype}
\usepackage{amsmath,amssymb,amsthm,mathtools}
\usepackage{siunitx}
\usepackage{graphicx}
\usepackage{booktabs}
\usepackage{array}
\usepackage[dvipsnames]{xcolor}
\usepackage{enumitem}
\usepackage[most]{tcolorbox}
\usepackage{tikz}
\usepackage{pgfplots}
\pgfplotsset{compat=1.18}
\usepackage[colorlinks=true,linkcolor=MidnightBlue,citecolor=ForestGreen,urlcolor=RoyalBlue]{hyperref}
\usepackage[nameinlink,capitalise]{cleveref}

\setlist[itemize]{topsep=0.35em,itemsep=0.15em}
\setlist[enumerate]{topsep=0.35em,itemsep=0.15em}
\setlength{\parskip}{0.45em}
\setlength{\parindent}{0pt}

\theoremstyle{definition}
\newtheorem{definition}{Definition}[section]
\newtheorem{example}{Example}[section]
\newtheorem{exercise}{Exercise}[section]

\theoremstyle{plain}
\newtheorem{proposition}{Proposition}[section]
\newtheorem{theorem}{Theorem}[section]

\newtcolorbox{intuitionbox}{
  colback=blue!3,
  colframe=blue!45!black,
  title=Intuition,
  sharp corners,
  boxrule=0.5pt
}

\newtcolorbox{warningbox}{
  colback=orange!5,
  colframe=orange!70!black,
  title=Warning,
  sharp corners,
  boxrule=0.5pt
}

\newtcolorbox{resultbox}{
  colback=green!4,
  colframe=green!45!black,
  title=Key result,
  sharp corners,
  boxrule=0.5pt
}

\newtcolorbox{summarybox}{
  colback=gray!5,
  colframe=gray!60!black,
  title=Summary,
  sharp corners,
  boxrule=0.5pt
}

\title{[Title from blueprint]}
\author{Educational note based on a YouTube transcript blueprint}
\date{\today}

\begin{document}

\maketitle

\begin{abstract}
[Brief summary of what the note teaches.]
\end{abstract}

\tableofcontents
\newpage

\section{[First section]}

[Introduce motivation, prerequisites, and notation before using formulas.]

\section{[Core development]}

\begin{definition}[Key term]
[Definition.]
\end{definition}

\begin{intuitionbox}
[Intuition connecting the formal idea to the learner's mental model.]
\end{intuitionbox}

\subsection{Step-by-step derivation}

State the goal before the equation block.
\begin{align}
  [starting point]
    &= [intermediate step] \\
    &= [next step] \\
    &= [final result].
\end{align}

Explain the role of each transformation in prose.

\begin{resultbox}
[Final result and interpretation.]
\end{resultbox}

\section{Worked examples}

\begin{example}[Example title]
[Step-by-step solution.]
\end{example}

\section{Conclusion}

[Synthesize the main result and what the learner should retain.]

\end{document}
